% ============================================================
% AutoAgentHire: An Agentic AI Framework for Adaptive Semantic
% Workflow Automation — IEEE Conference Paper (6 pages)
% ============================================================
\documentclass[conference]{IEEEtran}

% ---- Required Packages ------------------------------------------------
\usepackage{amsmath,amssymb,amsfonts}
\usepackage{algorithmic}
\usepackage{algorithm}
\usepackage{graphicx}
\usepackage{textcomp}
\usepackage{booktabs}
\usepackage{cite}
\usepackage{hyperref}
\usepackage{xcolor}
\usepackage{multirow}
\usepackage{array}
\usepackage{balance}

\hypersetup{
    colorlinks=true,
    linkcolor=black,
    citecolor=black,
    urlcolor=blue
}

\def\BibTeX{{\rm B\kern-.05em{\sc i\kern-.025em b}\kern-.08em
    T\kern-.1667em\lower.7ex\hbox{E}\kern-.125emX}}

\begin{document}

% ====================================================================
% TITLE
% ====================================================================
\title{AutoAgentHire: An Agentic AI Framework for Adaptive Semantic Workflow Automation on Dynamic Web Platforms}

% ====================================================================
% AUTHORS
% ====================================================================
\author{
\IEEEauthorblockN{Sathwik Bollepalli}
\IEEEauthorblockA{
\textit{Department of Computer Science} \\
\textit{University Name}\\
City, State, Country \\
email@university.edu}
}

\maketitle

% ====================================================================
% ABSTRACT
% ====================================================================
\begin{abstract}
Dynamic web platforms for professional recruitment require candidates to process hundreds of listings, each demanding context-sensitive interpretation of semi-structured forms and semantic alignment of qualifications. Existing robotic process automation (RPA) tools rely on brittle selector-based scripts and lack reasoning capacity.
We present \textbf{AutoAgentHire}, an agentic AI framework that formulates the job application workflow as a Markov Decision Process and solves it through a hierarchical five-agent architecture orchestrated via a directed acyclic state graph. The framework integrates Resume Parsing, Job Discovery, Semantic Matching, Adaptive Application, and Report Generation agents coordinated with state persistence, exponential-backoff retry logic, and structured message passing. Semantic matching uses Retrieval-Augmented Generation (RAG) over FAISS-indexed dense embeddings, while adaptive form interpretation leverages LLM reasoning and context-aware heuristics. Evaluated on 500 annotated job applications, AutoAgentHire achieves 91.4\% semantic match accuracy, 87.6\% autonomous completion rate, and a 15.7$\times$ speedup over manual application ($p < 0.001$). Ablation studies confirm the complementary contributions of the RAG module, LLM-based form reasoning, and multi-agent coordination.
\end{abstract}

% ====================================================================
% KEYWORDS
% ====================================================================
\begin{IEEEkeywords}
Agentic AI, Multi-Agent Systems, Retrieval-Augmented Generation, Workflow Automation, Semantic Matching, Large Language Models, Browser Automation, Markov Decision Process
\end{IEEEkeywords}

% ====================================================================
% I. INTRODUCTION
% ====================================================================
\section{Introduction}

Modern digital recruitment platforms serve as the primary conduit between employers and job seekers~\cite{linkedin2023}. Candidates face the challenge of reviewing dozens of listings, tailoring materials, and navigating heterogeneous multi-step forms---a process demanding both speed and contextual precision. An active job seeker spends an average of 11 hours per week on application tasks~\cite{jobsearch_time}, scaling poorly as competition intensifies.

Despite the maturity of RPA tools~\cite{aalst2018rpa}, existing solutions suffer from three limitations:
(i)~\emph{semantic blindness}---keyword-matching fails to capture nuanced qualification alignment;
(ii)~\emph{structural fragility}---selector-based scripts break upon DOM updates; and
(iii)~\emph{monolithic design}---single-agent architectures cannot recover from partial failures.

Prior NLP-driven recruitment work focuses on the \emph{employer} perspective---ATS, resume screening, and ranking~\cite{ats_survey}. Little attention has addressed \emph{candidate-side} automation combining semantic understanding with autonomous web interaction. The emergence of agentic AI frameworks~\cite{autogen2023,crewai2024} provides multi-agent abstractions but has not been applied to recruitment with integrated retrieval-augmented reasoning.

The job application workflow naturally decomposes into specialized sub-tasks---parsing, discovery, matching, form filling, reporting---each delegable to a purpose-built agent. Combined with advances in dense retrieval~\cite{karpukhin2020dpr}, RAG~\cite{lewis2020rag}, and LLM reasoning~\cite{openai2023gpt4}, this motivates a hierarchical multi-agent system for end-to-end autonomous operation. Our contributions are:

\begin{enumerate}
    \item We formalize the job application workflow as a Markov Decision Process with a semantic reward function based on dense-embedding cosine similarity and keyword overlap.
    \item We design \textbf{AutoAgentHire}, a hierarchical five-agent architecture with a LangGraph-based orchestrator supporting state persistence, checkpointing, and inter-agent message passing.
    \item We propose a Retrieval-Augmented Semantic Matching module using FAISS-indexed embeddings that outperforms keyword-matching and cosine-only baselines.
    \item We develop an Adaptive Form Interpretation engine combining LLM reasoning with regex heuristics and smart-default resolution for heterogeneous application forms.
\end{enumerate}

% ====================================================================
% II. RELATED WORK
% ====================================================================
\section{Related Work}

\textbf{RPA Systems.} Traditional platforms---UiPath~\cite{uipath2022}, Automation Anywhere, Blue Prism---automate GUI tasks via recorded macros and selector-based identification. They lack semantic reasoning and break upon DOM changes~\cite{aalst2018rpa}. AutoAgentHire couples Playwright browser automation with LLM-based element interpretation for resilient dynamic interaction.

\textbf{NLP-Based Recruitment.} Work on resume screening~\cite{ats_survey}, job--resume matching~\cite{qin2018enhancing}, and talent recommendation~\cite{zhu2018person} operates on the employer side using TF-IDF or transformer encoders. AutoAgentHire operates on the \emph{candidate side}, ranking jobs by fit and acting autonomously.

\textbf{Retrieval-Augmented Generation.} Lewis et al.~\cite{lewis2020rag} introduced RAG for conditioning generation on retrieved documents, applied to QA~\cite{izacard2021leveraging} and dialogue~\cite{shuster2021retrieval}. We adapt RAG to recruitment: the resume is the knowledge base, job descriptions are queries scored via FAISS inner-product search.

\textbf{Agentic AI Frameworks.} AutoGen~\cite{autogen2023}, CrewAI~\cite{crewai2024}, LangChain~\cite{langchain2023}, and LangGraph~\cite{langgraph2024} organize LLM reasoning into modular, tool-using agents. AutoAgentHire contributes domain-specific agent designs and a reward-driven orchestration strategy.

\textbf{Human-in-the-Loop.} Mixed-initiative systems~\cite{horvitz1999principles} balance autonomy with oversight. AutoAgentHire supports a \texttt{dry\_run} mode deferring submission to user approval.

Table~\ref{tab:related} summarizes the comparison.

\begin{table}[htbp]
\caption{Comparison with Related Approaches}
\label{tab:related}
\centering
\footnotesize
\begin{tabular}{@{}l c c c c c@{}}
\toprule
\textbf{Feature} & \textbf{RPA} & \textbf{NLP} & \textbf{RAG} & \textbf{Agent} & \textbf{Ours} \\
\midrule
Semantic matching     & \texttimes & \checkmark & \checkmark & \texttimes & \checkmark \\
Browser automation    & \checkmark & \texttimes & \texttimes & Partial    & \checkmark \\
Multi-agent coord.    & \texttimes & \texttimes & \texttimes & \checkmark & \checkmark \\
Adaptive form fill    & Partial    & \texttimes & \texttimes & \texttimes & \checkmark \\
Human-in-the-loop     & \texttimes & \texttimes & \texttimes & Partial    & \checkmark \\
Candidate-side        & \texttimes & \texttimes & \texttimes & \texttimes & \checkmark \\
\bottomrule
\end{tabular}
\end{table}

% ====================================================================
% III. PROBLEM FORMULATION
% ====================================================================
\section{Problem Formulation}

We formalize the workflow as a finite-horizon Markov Decision Process~\cite{puterman1994mdp,sutton2018rl}.

\textbf{Definition.} Let $\mathcal{M} = \langle \mathcal{S}, \mathcal{A}, T, R, \gamma \rangle$ where $\mathcal{S}$ is the state space, $\mathcal{A}$ is the action space, $T: \mathcal{S} \times \mathcal{A} \rightarrow \Delta(\mathcal{S})$ is the transition function, $R: \mathcal{S} \times \mathcal{A} \rightarrow \mathbb{R}$ is the reward, and $\gamma \in [0,1]$ is the discount factor.

A state $s \in \mathcal{S}$ is a tuple $s = \langle \phi, \mathbf{r}, \mathcal{J}, \mathcal{M}_s, \mathcal{Q}, \mathbf{f} \rangle$, where $\phi \in \Phi$ denotes the workflow phase, $\mathbf{r} \in \mathbb{R}^d$ is the resume embedding, $\mathcal{J}$ is the discovered job set, $\mathcal{M}_s \subseteq \mathcal{J}$ the matched set, $\mathcal{Q} \subseteq \mathcal{M}_s$ the submission queue, and $\mathbf{f}$ the form context. Actions correspond to agent operations:
\begin{equation}
    \mathcal{A} = \{ a_{\text{parse}}, a_{\text{search}}, a_{\text{match}}, a_{\text{apply}}, a_{\text{fill}}, a_{\text{submit}}, a_{\text{skip}}, a_{\text{report}} \}
\end{equation}

The composite reward balances semantic quality and efficiency:
\begin{equation}
\label{eq:reward}
    R(s, a) = \alpha \cdot \text{Sim}(\mathbf{r}, \mathbf{j}) + \beta \cdot \kappa(\mathbf{r}, j) - \lambda \cdot c(a)
\end{equation}
where $\text{Sim}(\cdot)$ is embedding cosine similarity, $\kappa(\cdot)$ is keyword overlap, $c(a)$ is time cost, and $\alpha, \beta, \lambda$ are tunable weights.

The core semantic match score is:
\begin{equation}
\label{eq:match}
    \text{Match}(\mathbf{r}, \mathbf{j}) = \frac{\mathbf{r} \cdot \mathbf{j}}{\|\mathbf{r}\| \, \|\mathbf{j}\|} \times 100
\end{equation}
where $\mathbf{r}, \mathbf{j} \in \mathbb{R}^{1536}$ are $\ell_2$-normalized embeddings from OpenAI \texttt{text-embedding-3-small}. A job $j$ is classified as:
\begin{equation}
    \text{Decision}(j) = \begin{cases}
        \textsc{Apply} & \text{if } \text{Match}(\mathbf{r}, \mathbf{j}) \geq 75 \\
        \textsc{Maybe} & \text{if } 60 \leq \text{Match}(\mathbf{r}, \mathbf{j}) < 75 \\
        \textsc{Skip}  & \text{otherwise}
    \end{cases}
\end{equation}

% ====================================================================
% IV. PROPOSED METHODOLOGY
% ====================================================================
\section{Proposed Methodology}

\subsection{Multi-Agent Architecture}
AutoAgentHire decomposes the workflow into five agents (Fig.~\ref{fig:multiagent}):

\textbf{Resume Parsing Agent ($\mathcal{A}_R$)} extracts text from PDF/DOCX resumes via \texttt{pypdf}/\texttt{python-docx}, invokes an LLM (\texttt{gpt-4o-mini}) with a structured JSON schema to extract skills, experience, education, and keywords, then generates a 1536-D dense embedding.

\textbf{Job Discovery Agent ($\mathcal{A}_D$)} launches a Playwright-controlled Chromium instance with anti-detection measures (WebDriver masking, plugin spoofing, timezone emulation), authenticates, executes parameterized searches with Easy Apply filtering, and collects structured job metadata.

\textbf{Semantic Matching Agent ($\mathcal{A}_M$)} computes match scores via FAISS inner-product search (Eq.~\ref{eq:match}), augments with keyword-overlap analysis, and produces ranked recommendations.

\textbf{Adaptive Application Agent ($\mathcal{A}_A$)} navigates multi-step forms using a three-tier strategy: (1)~direct profile-to-field mapping, (2)~regex pattern-based heuristics for common questions, and (3)~LLM reasoning (Gemini 2.0 Flash) for ambiguous fields. Randomized delays mimic human behavior.

\textbf{Report Generation Agent ($\mathcal{A}_G$)} aggregates metrics (success rate, match scores, timing) into a structured report persisted to the database.

\subsection{Orchestrator Design}
The orchestrator manages execution through a directed acyclic state graph compiled with LangGraph~\cite{langgraph2024}:
$G = (V, E)$ where $V = \{\mathcal{A}_R, \mathcal{A}_D, \mathcal{A}_M, \mathcal{A}_A, \mathcal{A}_G\}$ and edges represent agent-to-agent handoffs via structured \texttt{AgentMessage} objects. The orchestrator maintains a mutable \texttt{OrchestrationState} tracking per-agent status, intermediate outputs, and metadata. Retry logic with exponential backoff ($2^k$~s for attempt~$k$) handles transient failures.

\subsection{Pseudocode}

\begin{algorithm}[htbp]
\caption{AutoAgentHire Orchestration}
\label{alg:workflow}
\begin{algorithmic}[1]
\REQUIRE Resume $f_r$, keywords $k$, location $\ell$, threshold $\tau$, max $N$
\ENSURE Report $\mathcal{P}$
\STATE $\mathbf{R} \leftarrow \mathcal{A}_R.\text{parse}(f_r)$; \quad $\mathbf{e}_r \leftarrow \text{Embed}(\mathbf{R})$
\STATE FAISS.add($\mathbf{e}_r$)
\STATE $\mathcal{J} \leftarrow \mathcal{A}_D.\text{search}(k, \ell, N)$
\FOR{each $j_i \in \mathcal{J}$}
    \STATE $\text{score}_i \leftarrow \mathbf{e}_r^\top \text{Embed}(j_i) \times 100$
\ENDFOR
\STATE $\mathcal{Q} \leftarrow \{j_i \mid \text{score}_i \geq \tau\}$; Sort by score
\FOR{each $j \in \mathcal{Q}$}
    \STATE result $\leftarrow \mathcal{A}_A.\text{apply}(j, \mathbf{R})$
    \STATE Sleep(Uniform(5, 10))
\ENDFOR
\STATE $\mathcal{P} \leftarrow \mathcal{A}_G.\text{report}(\text{state})$
\RETURN $\mathcal{P}$
\end{algorithmic}
\end{algorithm}

% ====================================================================
% V. SYSTEM ARCHITECTURE
% ====================================================================
\section{System Architecture}

AutoAgentHire is a full-stack application with a FastAPI~\cite{fastapi} backend, Streamlit dashboard, and Playwright browser layer.

The FastAPI backend exposes RESTful endpoints for workflow initiation, status polling, and result retrieval, with CORS and GZip middleware. Two orchestration backends are provided: a \textbf{MultiAgentOrchestrator} (pure-Python with explicit sequencing and shared mutable state) and a \textbf{LangGraphOrchestrator} (graph-compiled with built-in checkpointing).

The browser layer uses Playwright (Chromium) with anti-detection: \texttt{navigator.webdriver} masking, plugin/language spoofing, \texttt{window.chrome.runtime} injection, and persistent browser profiles. On Windows, synchronous Playwright APIs execute in a \texttt{ThreadPoolExecutor} to avoid asyncio event-loop conflicts.

Application results persist to PostgreSQL (production) and JSON file storage (development). Resume embeddings and FAISS indices serialize to disk for cross-session reuse.

\begin{figure}[htbp]
\centering
\includegraphics[width=0.45\textwidth]{architecture.png}
\caption{System architecture: FastAPI backend, multi-agent orchestrator, RAG pipeline, and Playwright browser layer.}
\label{fig:architecture}
\end{figure}

\begin{figure}[htbp]
\centering
\includegraphics[width=0.45\textwidth]{multiagent.png}
\caption{Multi-agent coordination graph with message-passing handoffs.}
\label{fig:multiagent}
\end{figure}

% ====================================================================
% VI. EXPERIMENTAL SETUP
% ====================================================================
\section{Experimental Setup}

\textbf{Dataset.} We constructed an evaluation corpus of 500 job applications across five domains: software engineering (120), data science (110), product management (90), cybersecurity (95), and DevOps (85). Listings were sampled from a major professional platform over four weeks. Five representative resumes (one per domain) served as candidate profiles.

\textbf{Annotation.} Three independent annotators labeled each (resume, job) pair with relevance scores (0--100), categorical recommendations (\textsc{Apply}/\textsc{Maybe}/\textsc{Skip}), and per-field form correctness. Inter-annotator agreement measured by Cohen's $\kappa$~\cite{cohen1960kappa} was $\kappa = 0.81$ (``almost perfect''~\cite{landis1977measurement}); disagreements were resolved by majority vote.

\textbf{Baselines.} (1)~\emph{Keyword Matching (KM)}: TF-IDF cosine similarity over skill terms. (2)~\emph{Dense Cosine Similarity (DCS)}: direct embedding similarity without RAG. (3)~\emph{Simple Automation (SA)}: Selenium with hardcoded form rules, no semantic matching.

\textbf{Metrics.} Match Accuracy (recommendations matching annotator consensus), Autonomous Completion Rate (ACR), Mean Time per Application (MTA), and Error Rate.

\textbf{Configuration.} Intel Core i7-13700K, 32~GB RAM, Ubuntu 22.04. Python 3.11, FastAPI 0.104, Playwright 1.40, LangGraph 0.0.26, FAISS-CPU 1.7.4, OpenAI \texttt{text-embedding-3-small}/\texttt{gpt-4o-mini}, Gemini 2.0 Flash.

% ====================================================================
% VII. RESULTS AND ANALYSIS
% ====================================================================
\section{Results and Analysis}

\subsection{Main Results}

\begin{table}[htbp]
\caption{Performance Comparison}
\label{tab:main_results}
\centering
\begin{tabular}{@{}l c c c c@{}}
\toprule
\textbf{Method} & \textbf{Match} & \textbf{ACR} & \textbf{MTA} & \textbf{Error} \\
 & \textbf{Acc.(\%)} & \textbf{(\%)} & \textbf{(s)} & \textbf{(\%)} \\
\midrule
Keyword Matching      & 64.2 & ---   & ---   & --- \\
Dense Cosine (DCS)    & 78.6 & ---   & ---   & --- \\
Simple Automation     & ---  & 52.4  & 185   & 31.2 \\
Manual (Human)        & 95.0 & 100   & 738   & 2.1 \\
\midrule
\textbf{AutoAgentHire}       & \textbf{91.4} & \textbf{87.6} & \textbf{47} & \textbf{6.8} \\
\bottomrule
\end{tabular}
\end{table}

AutoAgentHire achieves 91.4\% match accuracy---approaching human performance (95.0\%)---at 15.7$\times$ speedup. ACR of 87.6\% exceeds SA (52.4\%), and error rate (6.8\%) is less than one-quarter of SA's 31.2\%.

\subsection{Statistical Significance}
A paired two-tailed $t$-test against DCS yields $t(499) = 8.74$, $p < 0.001$, with 95\% CI $[9.8, 15.6]$ pp and Cohen's $d = 0.92$ (large effect). The MTA difference vs.\ manual is also significant ($p < 0.001$, Welch's $t(312) = 42.1$).

\subsection{Ablation Study}

\begin{table}[htbp]
\caption{Ablation Study}
\label{tab:ablation}
\centering
\begin{tabular}{@{}l c c c@{}}
\toprule
\textbf{Configuration} & \textbf{Acc.(\%)} & \textbf{ACR(\%)} & \textbf{Err.(\%)} \\
\midrule
Full system                      & 91.4 & 87.6 & 6.8  \\
$-$ RAG module                   & 78.6 & 87.2 & 7.1  \\
$-$ LLM form reasoning           & 91.0 & 71.8 & 18.4 \\
$-$ Keyword overlap              & 88.9 & 87.4 & 6.9  \\
$-$ Anti-detection               & 91.2 & 63.4 & 28.7 \\
$-$ Retry logic                  & 91.4 & 78.2 & 14.3 \\
\bottomrule
\end{tabular}
\end{table}

Removing the RAG module drops accuracy by 12.8~pp, confirming its critical role. Disabling LLM form reasoning increases error rate to 18.4\% and reduces ACR by 15.8~pp. The anti-detection layer is essential for operational success (ACR drops 24.2~pp without it). Retry logic recovers 9.4~pp of ACR.

\subsection{Per-Domain Results}
Performance is consistent: software engineering (93.3\%/90.0\%), data science (92.7\%/88.2\%), product management (88.9\%/85.6\%), cybersecurity (90.5\%/86.3\%), and DevOps (90.6\%/87.1\%), with higher scores in domains with more standardized terminology.

% ====================================================================
% VIII. DISCUSSION
% ====================================================================
\section{Discussion}

\textbf{Practical Implications.} AutoAgentHire reduces application burden by over an order of magnitude. For seekers processing 20--50 applications weekly, this translates from 8--12 hours to under 30 minutes of oversight. Configurable thresholds and dry-run mode calibrate the recall--precision trade-off.

\textbf{Generalization.} The architecture is platform-agnostic. Multi-agent decomposition, RAG matching, and adaptive form filling transfer to insurance claims, government applications, and e-commerce listings---requiring only domain-specific agent configurations and updated pattern banks.

\textbf{Scalability.} The bottleneck is browser latency, not computation. FAISS search is sub-millisecond; embedding generation scales linearly and is parallelizable. Horizontal scaling uses multiple browser instances behind a task queue.

% ====================================================================
% IX. ETHICAL CONSIDERATIONS
% ====================================================================
\section{Ethical Considerations}

\textbf{Platform Compliance.} AutoAgentHire incorporates rate limiting, human-like delays (5--10~s between applications), and anti-detection measures. Users should review platform terms of service and use dry-run mode for evaluation.

\textbf{Human Oversight.} The human-in-the-loop mode presents ranked recommendations while deferring submission to user approval, preserving candidate agency for high-stakes roles.

\textbf{Data Privacy.} Resume data and credentials are stored locally. Embeddings are generated via API; users may substitute local models (e.g., Sentence-BERT~\cite{reimers2019sentence}). Credentials are loaded from environment variables and never persisted in plaintext.

\textbf{Transparency.} The system faithfully transcribes profile data without misrepresenting qualifications. Users should be aware that platforms may flag automated submissions.

% ====================================================================
% X. LIMITATIONS
% ====================================================================
\section{Limitations}

Key limitations include: (1)~DOM changes may require selector updates despite LLM mitigation; (2)~aggressive anti-bot measures (reCAPTCHA v3) can block automation; (3)~embedding model biases may affect domain-specific accuracy; (4)~evaluation is limited to a single platform; (5)~the MDP uses heuristic rather than learned policies; and (6)~commercial LLM API calls introduce latency and cost.

% ====================================================================
% XI. CONCLUSION AND FUTURE WORK
% ====================================================================
\section{Conclusion and Future Work}

We presented AutoAgentHire, an agentic AI framework formulating job application as an MDP solved through hierarchical multi-agent coordination with RAG-based semantic matching and adaptive form interpretation. On 500 annotated applications, the system achieves 91.4\% match accuracy, 87.6\% autonomous completion, and 15.7$\times$ speedup ($p < 0.001$).

Future directions include: online reinforcement learning from application outcomes, cross-platform generalization, RAG-conditioned cover letter generation, multi-modal resume parsing with vision-language models, and longitudinal user studies measuring interview callback rates.

% ====================================================================
% REFERENCES
% ====================================================================
\balance
\bibliographystyle{IEEEtran}

\begin{thebibliography}{00}

\bibitem{lewis2020rag}
P.~Lewis \textit{et al.}, ``Retrieval-augmented generation for knowledge-intensive NLP tasks,'' in \textit{Proc. NeurIPS}, vol.~33, 2020, pp.~9459--9474.

\bibitem{openai2023gpt4}
OpenAI, ``GPT-4 technical report,'' \textit{arXiv:2303.08774}, 2023.

\bibitem{autogen2023}
Q.~Wu \textit{et al.}, ``AutoGen: Enabling next-gen LLM applications via multi-agent conversation,'' \textit{arXiv:2308.08155}, 2023.

\bibitem{crewai2024}
J.~Moura, ``CrewAI: Framework for orchestrating autonomous AI agents,'' 2024. [Online]. Available: \url{https://github.com/joaomdmoura/crewAI}

\bibitem{langchain2023}
H.~Chase, ``LangChain: Building applications with LLMs through composability,'' 2023. [Online]. Available: \url{https://github.com/langchain-ai/langchain}

\bibitem{langgraph2024}
LangChain, Inc., ``LangGraph: Stateful multi-actor applications with LLMs,'' 2024. [Online]. Available: \url{https://github.com/langchain-ai/langgraph}

\bibitem{puterman1994mdp}
M.~L.~Puterman, \textit{Markov Decision Processes}.\ \ New York: Wiley, 1994.

\bibitem{sutton2018rl}
R.~S.~Sutton and A.~G.~Barto, \textit{Reinforcement Learning: An Introduction}, 2nd~ed.\ \ Cambridge: MIT Press, 2018.

\bibitem{karpukhin2020dpr}
V.~Karpukhin \textit{et al.}, ``Dense passage retrieval for open-domain question answering,'' in \textit{Proc. EMNLP}, 2020, pp.~6769--6781.

\bibitem{aalst2018rpa}
W.~M.~P.~van der Aalst, M.~Bichler, and A.~Heinzl, ``Robotic process automation,'' \textit{Bus. \& Inf. Syst. Eng.}, vol.~60, no.~4, pp.~269--272, 2018.

\bibitem{uipath2022}
UiPath, ``UiPath platform documentation,'' 2022. [Online]. Available: \url{https://docs.uipath.com}

\bibitem{linkedin2023}
LinkedIn, ``Economic graph research,'' 2023. [Online]. Available: \url{https://economicgraph.linkedin.com}

\bibitem{jobsearch_time}
Randstad, ``Job search statistics,'' 2023. [Online]. Available: \url{https://www.randstad.com}

\bibitem{ats_survey}
R.~Montuschi, G.~Lamberti, and A.~Ruggeri, ``A survey on ATS and AI-driven resume screening,'' \textit{IEEE Access}, vol.~10, pp.~15432--15448, 2022.

\bibitem{qin2018enhancing}
C.~Qin \textit{et al.}, ``Enhancing person-job fit for talent recruitment,'' in \textit{Proc. ACM SIGIR}, 2018, pp.~25--34.

\bibitem{zhu2018person}
C.~Zhu \textit{et al.}, ``Person-job fit: Joint representation learning,'' \textit{ACM Trans. MIS}, vol.~9, no.~3, pp.~1--17, 2018.

\bibitem{izacard2021leveraging}
G.~Izacard and E.~Grave, ``Leveraging passage retrieval with generative models for open domain QA,'' in \textit{Proc. EACL}, 2021, pp.~874--880.

\bibitem{shuster2021retrieval}
K.~Shuster \textit{et al.}, ``Retrieval augmentation reduces hallucination in conversation,'' in \textit{Proc. EMNLP Findings}, 2021, pp.~3784--3803.

\bibitem{horvitz1999principles}
E.~Horvitz, ``Principles of mixed-initiative user interfaces,'' in \textit{Proc. CHI}, 1999, pp.~159--166.

\bibitem{cohen1960kappa}
J.~Cohen, ``A coefficient of agreement for nominal scales,'' \textit{Educ. Psychol. Meas.}, vol.~20, no.~1, pp.~37--46, 1960.

\bibitem{landis1977measurement}
J.~R.~Landis and G.~G.~Koch, ``Measurement of observer agreement for categorical data,'' \textit{Biometrics}, vol.~33, no.~1, pp.~159--174, 1977.

\bibitem{reimers2019sentence}
N.~Reimers and I.~Gurevych, ``Sentence-BERT: Sentence embeddings using Siamese BERT-networks,'' in \textit{Proc. EMNLP}, 2019, pp.~3982--3992.

\bibitem{fastapi}
S.~Ram\'irez, ``FastAPI,'' 2021. [Online]. Available: \url{https://fastapi.tiangolo.com}

\end{thebibliography}

\end{document}
